\section{Agent-Based Model Equations}
\label{sec:abm_equations}

The crop-soil system is modeled as a set of $N=130$ coupled discrete-time dynamical systems. Each agent $n \in \{1, \dots, N\}$ maintains five state variables that evolve at a daily time step $k$: volumetric soil water content in the root zone ($x_1$, mm), cumulative thermal time ($x_2$, $^{\circ}$C$\cdot$day), cumulative maturation index ($x_3$), aboveground dry biomass ($x_4$, g/m$^2$), and surface ponding depth ($x_5$, mm).

To accurately simulate lateral water redistribution while maintaining computational efficiency for the optimal controller, the grid is solved sequentially at each time step $k$ following a topological sort of the directed elevation graph, executing from the highest elevation agents down to the sinks. Furthermore, horizontal subsurface lateral flow is considered negligible for the 460~m$^2$ agent resolution; thus, agents are coupled exclusively via surface runoff.

\subsection{Surface Hydrology and Runoff}

The total available surface water ($W_{\text{surf}}$) for agent $n$ consists of yesterday's un-infiltrated ponding ($x_5$), daily precipitation ($P$), applied irrigation ($u^n$), and incoming surface runoff routed from all uphill neighbors $m \in \mathcal{U}(n)$, divided equally among the neighbor's downhill receivers ($N_r^m$):
\begin{equation}
    W_{\text{surf}}^n(k) = x_5^n(k) + P(k) + u^n(k) + \sum_{m \in \mathcal{U}(n)} \frac{\varphi_2^m(k)}{N_r^m}
\end{equation}

Surface runoff ($\varphi_2$) is computed using the SCS curve number method. To enforce strict mass conservation and simulate topographical valley flooding, a boundary condition is applied where sink agents ($N_r^n = 0$) cannot generate runoff:
\begin{equation}
    \varphi_2^n(k) = 
    \begin{cases}  
      \frac{(W_{\text{surf}}^n(k) - \theta_3)^2}{W_{\text{surf}}^n(k) + 4\theta_3}, & N_r^n > 0 \text{ and } W_{\text{surf}}^n(k) > \theta_3 \\
      0, & \text{otherwise} 
   \end{cases}
\end{equation}

\subsection{Subsurface Hydrology and Infiltration}

Crop transpiration ($\varphi_1$) is the water extracted by plant roots, bounded by the specific crop evapotranspiration ($\text{ET}_c$), which is derived from the atmospheric reference demand ($\text{ET}_0$) and the crop coefficient ($K_c$):
\begin{equation}
    \text{ET}_c(k) = K_c \cdot \text{ET}_0(k)
\end{equation}
\begin{equation}
    \varphi_1^n(k) = \min\!\Big(\max\!\big(\theta_1(x_1^n(k) - \theta_2\theta_5),\; 0\big),\; \text{ET}_c(k)\Big)
\end{equation}

Water that does not run off the surface infiltrates into the root zone. To prevent non-physical infinite pooling within the soil matrix, infiltration ($I$) is bounded by the physical saturation capacity of the soil ($\theta_{\text{sat}}\theta_5$):
\begin{equation}
    I_{\max}^n(k) = \theta_{\text{sat}}\theta_5 - x_1^n(k) + \varphi_1^n(k)
\end{equation}
\begin{equation}
    I^n(k) = \min\!\big(W_{\text{surf}}^n(k) - \varphi_2^n(k),\; I_{\max}^n(k)\big)
\end{equation}

An intermediate soil moisture state ($x_{1,\text{temp}}$) is computed to evaluate the temporary water holding before daily drainage occurs:
\begin{equation}
    x_{1,\text{temp}}^n(k) = x_1^n(k) + I^n(k) - \varphi_1^n(k)
\end{equation}

If this temporary moisture exceeds the field capacity ($\theta_6\theta_5$), the subsurface excess ($E_{\text{sub}}$) becomes vulnerable to gravity. A fraction of this excess, defined by the daily drainage rate ($\theta_4$), is lost to deep vertical drainage ($\varphi_3$):
\begin{equation}
    E_{\text{sub}}^n(k) = \max\!\big(x_{1,\text{temp}}^n(k) - \theta_6\theta_5,\; 0\big)
\end{equation}
\begin{equation}
    \varphi_3^n(k) = \theta_4 \cdot E_{\text{sub}}^n(k)
\end{equation}

\subsection{Stress Penalties and Crop Growth}
\label{sec:stress_growth}

Thermal time accumulation ($h_1$) drives phenological development:
\begin{equation}
    h_1^n(k) = \max\!\big(T_{\text{mean}}(k) - \theta_7,\; 0\big)
\end{equation}

Heat stress ($h_2$) reduces maturation efficiency, and low-temperature stress ($h_7$) prevents biomass accumulation:
\begin{equation}
    h_2(k) = \text{clip}\!\left( \frac{\theta_{10} - T_{\max}(k)}{\theta_{10} - \theta_9},\; 0,\; 1 \right)
\end{equation}
\begin{equation}
    h_7(k) = 
    \begin{cases} 
      1, & T_{\text{mean}}(k) > \theta_7 \\ 
      0, & T_{\text{mean}}(k) \leq \theta_7 
   \end{cases}
\end{equation}

Hydrological stresses directly impact crop development. Drought stress ($h_3$) penalizes the system when actual transpiration falls below the crop's atmospheric demand ($\text{ET}_c$), while waterlogging stress ($h_6$) penalizes root-zone oxygen deprivation when soil moisture exceeds field capacity:
\begin{equation}
    h_3^n(k) = 1 - \theta_{14} \cdot \max\!\left(1 - \frac{\varphi_1^n(k)}{\text{ET}_c(k)},\; 0\right)
\end{equation}
\begin{equation}
    h_6^n(k) = \text{clip}\!\left(1 - \frac{x_1^n(k) - \theta_6\theta_5}{\theta_6\theta_5},\; 0,\; 1\right)
\end{equation}

The growth function ($g$) dictates the fraction of intercepted solar radiation contributing to biomass, forming a two-branch sigmoid curve scaled by the maximum interception parameter ($\theta_{19}$):
\begin{equation}
g(x_2^n) = \begin{cases}
\theta_{19}/(1 + e^{-0.01(x_2^n - \theta_{20})}) & x_2^n \leq \theta_{18}/2 \\
\theta_{19}/(1 + e^{0.01(x_2^n + x_3^n - \theta_{18})}) & x_2^n > \theta_{18}/2
\end{cases}
\end{equation}


\subsection{State Update Equations}

The coupled dynamics for agent $n$ conclude with the daily state updates. Any surface water that failed to infiltrate remains explicitly conserved in the surface ponding state ($x_5$) for the subsequent day:
\begin{align}
    x_1^n(k+1) &= x_{1,\text{temp}}^n(k) - \varphi_3^n(k) \\
    x_2^n(k+1) &= x_2^n(k) + h_1^n(k) \\
    x_3^n(k+1) &= x_3^n(k) + \theta_{11}\big(1 - h_2(k)\big) + \theta_{12}\big(1 - h_3^n(k)\big) \\
    x_4^n(k+1) &= x_4^n(k) + \theta_{13} \cdot h_3^n(k) \cdot h_6^n(k) \cdot h_7(k) \cdot g(x_2^n(k)) \cdot R_s(k) \\
    x_5^n(k+1) &= W_{\text{surf}}^n(k) - \varphi_2^n(k) - I^n(k)
\end{align}

\section{Model Corrections and Extensions}
\label{sec:corrections}

While the baseline framework~\cite{Lopezjimenez2024} provides a robust topographical architecture, implementation for closed-loop control revealed necessary mathematical corrections to enforce strict mass-balance and to ensure controller cost-functions were properly coupled to agronomic penalties. Five primary modifications were integrated into the environment.

\subsection{Hydrological Decoupling and Mass Conservation}

The original formulation in~\cite{Lopezjimenez2024} computes an outflow term ($\phi_{\text{out}}$) that lumps surface runoff and subsurface excess together. While the runoff portion is properly subtracted from the sender's state equation, the subsurface excess portion is not. Consequently, a sending agent passes its subsurface excess to downhill neighbors but never loses that volume from its own root zone. This generates an unbounded artificial creation of water ("phantom water"), leading to exaggerated flooding in lower elevations during sustained rainfall. 

The proposed decoupled two-layer architecture resolves these flaws. By explicitly bridging the surface layer ($W_{\text{surf}}$) and the root zone ($x_1$) via a strictly bounded infiltration term ($I$), the overlapping fluxes are eliminated. Explicitly tracking the exact volume of water moving between the decoupled layers guarantees strict mass conservation across the entire field.

\subsection{Surface Routing, Sinks, and Ponding Dynamics}
\label{sec:runoff}

In the original configuration, surface runoff ($\varphi_2$) is mathematically aggregated with subsurface flow and routed directly into the soil matrix of adjacent agents. This simplification bypasses surface pooling dynamics, injecting rapid overland flow directly into the root zone of downhill neighbors. 

To more accurately capture the bowl-shaped topography of the field, surface runoff is explicitly decoupled from subsurface flow in this work. Runoff generated by an agent is routed to lower-elevation neighbors where it arrives as surface inflow. This inflow is added to the neighbor's available surface water ($W_{\text{surf}}$) and must undergo infiltration before entering the root zone. 

Crucially, a topological boundary condition was introduced to handle valley flooding. In the baseline model, sink agents ($N_r = 0$) generated runoff that mathematically vanished. By enforcing $\varphi_2 = 0$ for all sinks, and by capturing un-infiltrated water in a new surface ponding state variable ($x_5$), the model realistically simulates sustained, multi-day waterlogging in low-elevation zones.

To justify the extra computational cost of routing water across the elevation map, the cascade model was compared against a simpler single-step baseline where each grid cell acts independently. Both models were tested under the wet year scenario (2024) with three different irrigation levels (0 mm, 5 mm, and 8 mm). While the simpler model produces similar field-average yields, the cascade mode reveals substantially different sink agent behavior (Figure~\ref{fig:runoff_comparison_grid}), particularly under irrigation. Because the MPC controller must predict and avoid waterlogging at sink agents, the cascade mode was adopted to ensure the prediction model accurately reflects within-day runoff accumulation.

\begin{figure}[H]
    \centering
    \includegraphics[width=1\textwidth]{figures/fig_runoff_comparison_grid.png}
    \caption{Hydrological comparison of runoff modes across 0, 5, and 8 mm/day irrigation levels (2024)}
    \label{fig:runoff_comparison_grid}
\end{figure}

\subsection{Scale Justification: Omission of Subsurface Lateral Flow}

The baseline model processes horizontal subsurface lateral flow at the same discrete daily time step as rapid surface runoff. However, given the 460~m$^2$ spatial resolution of the agents (approximately $21 \times 21$~m plots), lateral water movement through dense silty loam over a single day is physically negligible compared to vertical deep drainage and overland flow. 

Therefore, horizontal subsurface flow ($\varphi_4$) was omitted from the proposed framework. The root zones are modeled as independent vertical columns coupled exclusively by the surface runoff routing graph. This physical simplification significantly increases the sparsity of the Jacobian matrices, thereby reducing the computational complexity and accelerating the prediction horizon of the MPC.

\subsection{Agronomic Penalties and Crop Evapotranspiration}

In the original framework, transpiration and drought stress were bounded by the reference evapotranspiration ($\text{ET}_0$), effectively treating all crops equivalently to a standard 12~cm reference grass. This formulation was corrected by introducing crop coefficients ($K_c$) to calculate crop-specific evapotranspiration ($\text{ET}_c$), ensuring Hashemi rice and tobacco respect their distinct biological water demands.

Furthermore, in the original model, the drought stress function $h_3$ appears only in the maturation equation ($x_3$), where it accelerates crop senescence under water deficit. However, $h_3$ is absent from the biomass equation ($x_4$). This means that a completely unirrigated crop would accumulate the same daily biomass increment as a fully irrigated one, provided all other conditions were equal. The only effect of drought would be to advance the onset of senescence, which reduces biomass accumulation only indirectly and with considerable delay.

This omission was identified during the initial ABM validation, when the no-irrigation baseline produced biomass values that were nearly indistinguishable from irrigated scenarios for wheat. The corrected formulation multiplies the radiation use efficiency term ($\theta_{13}$) in the biomass state update by $h_3$, so that drought directly reduces the daily biomass increment. This modification follows the approach of the FAO AquaCrop model, in which water stress reduces radiation use efficiency proportionally to the ratio of actual to potential transpiration. With $\theta_{14} = 0.8$ for rice and 0.6 for tobacco, maximum drought reduces the biomass increment by up to 80\% and 60\%, making irrigation decisions directly consequential for yield outcomes.

This correction is particularly important for the controller comparison objective of this thesis. Without it, the ABM would produce only marginal differences between irrigated and unirrigated outcomes, rendering any comparison between MPC, RL, and baseline controllers meaningless. The original crop used in the baseline paper was wheat in Colombia, which is less water-sensitive (lower $\theta_{14}$); it is possible that the omission went undetected in that context because the irrigation response was already small.

\subsection{Biomass Calibration and Initialization}

The radiation use efficiency parameter $\theta_{13}$ was calibrated so that $x_4$ reliably represents g/m$^2$ of aboveground dry matter. The final agricultural yield is calculated as $Y = x_4 \times \text{HI} \times 10$ (kg/ha). 

For Hashemi rice, the simulated yields were validated against typical Gilan farmer reports ($\sim$2000~kg/ha) and optimal experimental conditions (3500--4000~kg/ha)~\cite{jalali2021}. The initial state conditions reflecting the 20-day nursery period for rice are set to $x_{2,\text{init}} = 210$ and $x_{4,\text{init}} = 60$. Conversely, tobacco seedlings are initialized at $x_{2,\text{init}} = 0$ and $x_{4,\text{init}} = 8$.

\section{Soil and Crop Parameterization}
\label{sec:soil_crop_params}

\subsection{Soil Parameters}
\label{soil_param}
The soil at the study site is classified as silty loam, based on regional soil surveys and FAO soil classification maps for Gilan Province. The soil hydraulic properties used in this study follow FAO Irrigation and Drainage Paper No.~56 recommendations for this soil type. 

\begin{table}[htbp]
\centering
\caption{Soil hydraulic parameters for silty loam (common to all crops)}
\label{tab:soil}

\begin{tabular}{cccc}
\toprule
\textbf{Parameter} & \textbf{Symbol} & \textbf{Value} & \textbf{Source}\\
\midrule
Water uptake coefficient & $\theta_1$ & 0.096 & \cite{Lopezjimenez2024} \\
Wilting point (volumetric fraction) & $\theta_2$ & 0.15 & \cite{allen1998} \\
SCS threshold (mm) & $\theta_3$ & 10.0 & \cite{Lopezjimenez2024} \\
Drainage coefficient & $\theta_4$ & 0.05 & \cite{Lopezjimenez2024} \\
Field capacity (volumetric fraction) & $\theta_6$ & 0.35 & \cite{allen1998} \\
Saturation (volumetric fraction) & $\theta_{sat}$ & 0.50 & \cite{rawls1982} \\
\bottomrule
\end{tabular}
\end{table}

\subsection{Crop Parameters}
\label{crop_param}

Two crops were selected for the simulation environment: rice (Hashemi cultivar) and tobacco (\emph{Nicotiana tabacum}). Both hold significant commercial value in Gilan Province~\cite{surfiran2025, iranica_tobacco}. Their contrasting agronomic characteristics (high water demand, shallow roots, and high drought sensitivity for rice, versus moderate water demand, deep roots, and relative drought tolerance for tobacco) provide a robust basis for evaluating the control algorithms under differing physical constraints.

\subsubsection{Rice (Hashemi Cultivar)}

The base temperature ($T_{\text{base}}$ or $\theta_7$) is set to $10^{\circ}$C, following standard Growing Degree Day (GDD) conventions~\cite{mcmaster1997, laohasiriwong2023} and the baseline ABM~\cite{Lopezjimenez2024}. The maturity threshold ($\theta_{18}$) was calibrated to 1250~$^{\circ}$C$\cdot$day based on the 25-year dataset and validated against local field measurements, which report 1216--1352~GDD for the Hashemi cultivar in Gilan~\cite{sadidi2021}. 

Transplanting is simulated on May 20 (DOY 141), aligning with local practices when the 7-day rolling average $T_{\text{mean}} \geq 18^{\circ}$C and $T_{\min} \geq 12^{\circ}$C~\cite{paredes2025, faghani2011}. The harvest is set for August 20 (DOY 233). Extending the season beyond this date sharply increases the risk of extreme WMO R20mm rainfall events ($\geq$20~mm/day), which can cause crop lodging and severe harvest disruption~\cite{sarma2018, nayak2011}. 

This yields a 93-day field season. As established in Section~\ref{sec:corrections}, the 20-day prior nursery period contributes an initial thermal accumulation of $x_{2,\text{init}} = 210$~GDD and an initial biomass of $x_{4,\text{init}} = 60$~g/m$^2$. Other key parameters include a crop coefficient ($K_c$) of 1.15, a shallow root depth ($\theta_5$) of 400~mm, and a high drought sensitivity ($\theta_{14}$) of 0.80.

\subsubsection{Tobacco (\emph{Nicotiana tabacum})}

The base temperature is similarly set to $10^{\circ}$C, which falls within the standard 10--13$^{\circ}$C restriction range for tobacco~\cite{li2021}. The maturity threshold ($\theta_{18}$) is calibrated to 1200~GDD. This duration corresponds to a 104-day field season, consistent with standard FAO timelines (90--120 days)~\cite{fao_tobacco} and general agronomic data (100--130 days)~\cite{britannica1999}. As tobacco seedlings are typically delivered directly from nurseries with unknown prior conditions, the initial states are set to $x_{2,\text{init}} = 0$ and a minimal $x_{4,\text{init}} = 8$~g/m$^2$.


Transplanting occurs on May 25 (DOY 146), with harvest on September 5 (DOY 249). Tobacco features a significantly deeper root zone ($\theta_5 = 700$~mm) compared to rice, allowing it to access a larger reservoir of deep soil moisture. It also exhibits a lower crop coefficient ($K_c = 0.90$) and lower drought sensitivity ($\theta_{14} = 0.60$), making it more resilient to constrained irrigation schedules. 

\subsubsection{Parameter Summary}

\begin{table}[H]
\centering
\caption{Soil and crop-specific parameters used in the ABM simulation}
\label{tab:params}
\small
\begin{tabular}{llccl}
\toprule
\textbf{Parameter} & \textbf{Symbol} & \textbf{Rice} & \textbf{Tobacco} & \textbf{Source} \\
\midrule
Root depth & $\theta_5$ & 400 mm & 700 mm & \cite{allen1998}, \cite{fao_tobacco} \\
Base temperature & $\theta_7$ & 10$^{\circ}$C & 10$^{\circ}$C & \cite{laohasiriwong2023}, \cite{li2021} \\
Heat stress onset & $\theta_9$ & 35$^{\circ}$C & 30$^{\circ}$C & \cite{fao_tobacco} \\
Extreme heat & $\theta_{10}$ & 42$^{\circ}$C & 38$^{\circ}$C & -- \\
Heat stress maturation rate & $\theta_{11}$ & 0.0030 & 0.0025 & \cite{Lopezjimenez2024} \\
Drought stress maturation rate & $\theta_{12}$ & 0.0030 & 0.0028 & \cite{Lopezjimenez2024} \\
RUE & $\theta_{13}$ & 0.65 & 0.30 & \cite{patel2013}, \cite{ferrara2002} \\
Drought sensitivity & $\theta_{14}$ & 0.80 & 0.60 & \cite{Lopezjimenez2024} \\
Maturity GDD & $\theta_{18}$ & 1250 & 1200 & \cite{sadidi2021}, calib. \\
Max interception & $\theta_{19}$ & 0.95 & 0.90 & \cite{Lopezjimenez2024} \\
50\% intercept. GDD & $\theta_{20}$ & 417 & 450 & \cite{zhang2017} \\
Crop coefficient & $K_c$ & 1.15 & 0.90 & \cite{allen1998}, \cite{fao_tobacco} \\
Depletion fraction & $p$ & 0.20 & 0.50 & \cite{allen1998}, \cite{fao_tobacco} \\
Harvest index & HI & 0.42 & 0.55 & \cite{jalali2021}, \cite{fao_tobacco} \\
Nursery GDD & $x_{2,\text{init}}$ & 210 & 0 & Calib. \\
Initial biomass & $x_{4,\text{init}}$ & 60 g/m$^2$ & 8 g/m$^2$ & Est. \\
\bottomrule
\end{tabular}
\end{table}